\makeatletter
\renewcommand{\abstract@toptext}{
\noindent Proposta de Dissertação de Mestrado apresentada à COPPE/UFRJ como parte dos requisitos necessários para a obtenção do título de Mestre em Engenharia Mecânica.
}
\makeatother
\begin{abstract}

%A destilação por membranas aplicada ao tratamento de água de produção apresenta-se como uma tecnologia promissora para recuperação de água em processos da indústria do petróleo, devido à sua elevada rejeição de sais e capacidade de operar com soluções de alta salinidade. Entretanto, a ocorrência de mecanismos de falha, como molhamento (wetting) e incrustação (fouling), compromete a estabilidade operacional do processo ao longo do tempo, tornando necessária a avaliação do desempenho temporal das membranas empregadas. Nesse contexto, este trabalho propõe a análise experimental de membranas comerciais submetidas ao tratamento de água de produção por destilação por membranas com espaçamento de ar (AGMD), visando identificar a influência das propriedades morfológicas e operacionais na estabilidade do fluxo de permeado e na qualidade da água produzida. Para isso, serão avaliadas membranas fabricadas com diferentes materiais e suportes, previamente caracterizadas por porosimetria de fluxo capilar, microscopia digital e medições de ângulo de contato, permitindo determinar parâmetros como distribuição de poros, porosidade e pressão de entrada do líquido (LEP). Os experimentos serão conduzidos em uma bancada laboratorial instrumentada, operando com água de produção real oriunda de um campo petrolífero brasileiro, considerando condições controladas de temperatura, vazão e salinidade durante ensaios de longa duração. Além disso, pretende-se correlacionar os mecanismos de degradação observados com as características físico-químicas da água de produção e os parâmetros operacionais do processo, contribuindo para o desenvolvimento de sistemas de destilação por membranas mais confiáveis e adequados para aplicações industriais.
  

A água de produção é o principal resíduo gerado na exploração e produção de petróleo, sendo composta por água de formação, água de injeção e compostos oriundos do reservatório. Sua composição pode incluir sais dissolvidos, hidrocarbonetos, sólidos suspensos, metais e produtos químicos utilizados no processo produtivo. As características desse resíduo variam conforme a localização do campo petrolífero, as condições geológicas do reservatório e o estágio de produção do poço, tornando seu tratamento um desafio ambiental e operacional. Neste contexto, a presente dissertação tem como objetivo avaliar o desempenho de longo prazo de membranas comerciais no tratamento de água de produção via destilação por membranas, utilizando a configuração AGMD. O estudo busca investigar a influência das características morfológicas das membranas sobre o desempenho operacional e a resistência ao molhamento. As etapas já desenvolvidas incluíram a caracterização das membranas comerciais, abrangendo distribuição e diâmetro médio de poros, ângulo de contato e LEP. Também foi realizada a caracterização da água de produção utilizada nos ensaios, determinando as concentrações dos contaminantes, além de experimentos de curto prazo para avaliação preliminar do fluxo de permeado, estabilidade operacional e suscetibilidade ao molhamento das membranas. Os experimentos de longo prazo serão conduzidos com monitoramento contínuo de fluxo de permeado, condutividade, pressão e temperatura. A partir desses ensaios, serão avaliadas a estabilidade operacional das membranas, a ocorrência de molhamento e a degradação de desempenho ao longo do tempo, buscando correlacionar as propriedades morfológicas das membranas com seu comportamento operacional em aplicações prolongadas.

\end{abstract}

%\makeatletter
%\renewcommand{\foreignabstract@toptext}{
%\noindent Proposal of Master Dissertation presented to COPPE/UFRJ as a partial fulfillment of the requirements for the degree of Master in Mechanical Engineering.
%}
%\makeatother

%\begin{foreignabstract}

%Produced water is the main waste generated during oil exploration and production, consisting of formation water, injection water, and compounds originating from the reservoir. Its composition may include dissolved salts, hydrocarbons, suspended solids, metals, and chemical additives used in the production process. The characteristics of this effluent vary according to the location of the oilfield, the geological conditions of the reservoir, and the production stage of the well, making its treatment both an environmental and operational challenge. In this context, the present dissertation aims to evaluate the long-term performance of commercial membranes in the treatment of produced water via membrane distillation using the AGMD configuration. The study seeks to investigate the influence of membrane morphological characteristics on operational performance and wetting resistance. The stages already completed included the characterization of commercial membranes, encompassing pore size distribution, mean pore diameter, contact angle, and LEP. The produced water used in the experiments was also characterized by determining contaminant concentrations, in addition to conducting short-term experiments for the preliminary evaluation of permeate flux, operational stability, and membrane wetting susceptibility. Long-term experiments will be carried out with continuous monitoring of permeate flux, conductivity, pressure, and temperature. Based on these tests, membrane operational stability, wetting occurrence, and performance degradation over time will be evaluated, aiming to correlate membrane morphological properties with their operational behavior under prolonged application conditions.
  
%\end{foreignabstract}